\section{The Collaboration Model}
\label{sec:collaboration}

\subsection{A Language That Emerged From Conversation}
\label{sec:conversation}

Trenza was not designed in the conventional sense. There was no prior
specification from which an implementation was derived. There was a
conversation---sustained over six weeks, across multiple sessions and multiple
models---and the language is what that conversation converged on.

This is not an incidental feature of the project's history. It is the central
claim of this section: the design process that produced Trenza \emph{is} an
instance of the kind of human-LLM collaboration that Trenza is designed to
support. The chronicle of the project---more than ninety dated and attributed
entries committed to version control---is not documentation of a requirements
engineering process. It is the requirements engineering process itself.

The distinction matters. Post-hoc documentation reconstructs decisions after
they have been made; it tends to smooth over the disagreements that produced
the final shape and to attribute decisions to whoever happened to be writing
the documentation. The chronicle records decisions as they are made, including
the reasoning, the alternatives considered, the disagreements and their
resolution, and the open questions that remain. When a design decision is later
challenged, the challenge can be located in the record. When a model resumes
work after an interruption---and every model interruption is total, since no
model session has memory of any prior session---it reads the chronicle rather
than reconstructing context from code. The repository functions as the
persistent memory that no individual model session possesses.

\subsection{The Division of Roles}
\label{sec:roles}

The collaboration protocol that governs the project formalizes a division of
roles that emerged empirically and was subsequently codified in project
documentation and the design-decision record:

\begin{quote}
\itshape
The human provides intent. The model crystallizes specification. The compiler
verifies.
\end{quote}

In practice, this means: the human architect identifies the behavioral property
that needs to be expressed---\emph{``a flag checked in four places should be a
compile-time error''}---and the model translates that intent into a formal
construct. The human does not write Trenza syntax. The model does not decide
what Trenza should express. The compiler does not interpret either.

This division is not incidental to intellectual property considerations, though
it has implications for them. It is a structural property of the collaboration:
the model's contribution is genuinely productive---it identifies constructs,
proposes grammar, argues for and against design choices---but the decisions are
always ratified by the human before they are committed. The commit history is
the record of ratification. An uncommitted design decision is a proposal; a
committed one is a fact.

The multi-model character of the collaboration produces a form of distributed
peer review that a single-model workflow cannot replicate. Three models with
distinct profiles participated:

\begin{itemize}
  \item \textbf{Claude Sonnet~4.6} acted as session coordinator and primary
        author of specification text. It maintained narrative continuity across
        sessions by reading and writing the chronicle.
  \item \textbf{Claude Opus~4.6} acted as architectural reviewer. Its sessions
        were shorter and focused on structural critique: cross-checking design
        decisions, detecting incoherence between proposed changes and prior
        commitments, and proposing rules for the verifier.
  \item \textbf{Gemini~3.1 Pro} was nominated as implementer of the Rust compiler
        and adversarial reviewer of the language design. In practice, quota
        limits on Pro meant that the majority of the implementation
        work---the validator code, the AST, and a substantial portion of the
        generators---was carried out by \textbf{Gemini~3 Flash}, the smaller
        model, with Pro reserved for design discussion and architectural review.
\end{itemize}

The clearest single instance of distributed peer review producing something the
human would not have specified is the addition of \textbf{Rule~8
(role-type consistency)} to the verifier. The rule was not requested. A
Gemini~3 Flash session---that is, the smaller of the two Gemini models in
routine use on the project, used primarily for high-throughput implementation
work rather than design---was working on unrelated self-hosting tasks. While
verifying that the CLI specification in \texttt{trenza-cli.trz} could be
compiled by the CLI itself, the model identified a class of inconsistency that
Rules~1--7 did not catch: a role declared with one event surface in one context
could appear with a divergent surface in another, and nothing flagged it. The
model proposed the rule, implemented it, and added the corresponding test---all
in the same session, without prompting from the human or from the larger
Gemini~3.1 Pro model.

The anecdote matters for two reasons beyond the rule itself. First, the
contribution came from the \emph{smaller} model. The kind of work that matters
for a formal specification language is not always the kind of work that requires
the most capable model; it requires the model that is in the right context with
the right artifacts to see the gap. Second, the rule the model added is a rule
about \emph{naming consistency across files}---exactly the class of bug that
motivated the language in the first place. The collaboration produced not just
an arbitrary new verification rule but the rule that closes the same kind of
dispersal that produced the original CronometroPSP defect, applied now to role
declarations rather than to flag variables. The system extended itself along
its own grain.

\subsection{The Specification as Epistemic Artefact}
\label{sec:epistemic}

The most concrete result of the collaboration model is the \texttt{.trz}
specification as a reviewable artefact. An experiment conducted on 23 March
2026 tested whether the presence of a \texttt{.trz} specification changes the
quality of LLM-assisted code review. Three deliberate bugs were injected into
a generated Rust file: a semantic violation (\texttt{forbidden} silently
replaced by \texttt{ignored}), a structural error (an authentication failure
transitioning to an active session), and a completeness gap (a missing
\texttt{\#[should\_panic]} test).

Two review passes were conducted: one with the Rust file alone, one with the
\texttt{.trz} specification as reference. Both passes detected all three bugs.
The difference was not in what was found but in how:

\begin{table}[h]
  \caption{LLM code review with and without a \texttt{.trz} specification.
           Both passes detected all three injected bugs; the regime of work
           differed.}
  \label{tab:review}
  \small
  \begin{tabular}{p{3.0cm}p{1.5cm}p{1.9cm}}
    \toprule
    Dimension & Without spec & With spec \\
    \midrule
    Total reasoning steps                            & $\sim$22  & $\sim$7 \\
    Confidence in Bug~1 (\texttt{forbidden}$\to$\texttt{ignored}) & Medium & High \\
    Confidence in exhaustiveness                     & Low       & High \\
    Nature of work                                   & Heuristic & Mechanical \\
    \bottomrule
  \end{tabular}
\end{table}

The key finding was not a reduction in review time. It was an epistemic regime
shift. Without a specification, the reviewer produces probabilistic conclusions:
\emph{``this looks like a bug''}. With a specification, it produces deductive
conclusions: \emph{``this is a bug''}. The distinction is not academic. Bug~1
---\texttt{forbidden} replaced by \texttt{ignored}---produces identical
observable runtime behavior in both cases. The difference exists only in the
system's design contract. Without the contract made explicit, a reviewer cannot
confirm with certainty whether a divergence is a bug or a design decision. The
specification collapses that ambiguity to a string comparison.

This regime shift matters most for the class of defects that Trenza was designed
to prevent: semantic violations with no visible runtime signature, completeness
gaps that are undetectable without knowing the full intended contract. The
\texttt{.trz} is not documentation that accompanies the code; it is the ground
truth against which the code is verified. That the ground truth is
machine-readable---that it can be traversed exhaustively by a model, not just
cited selectively---is what makes the verification deductive rather than
heuristic.

\subsection{The Human as Architect}
\label{sec:architect}

The collaboration model described above presupposes a particular kind of human
participation. The human is not a ``code typist'' who specifies implementation
details. Nor is the human a passive approver who accepts whatever the model
generates. The human is an architect in the original sense: one who determines
what the system must be, and delegates the question of how it is expressed to a
disciplined process.

This role requires a specific kind of attention: not to syntactic details, but
to behavioral contracts. The question is not \emph{``does this code look
right?''} but \emph{``does the system's behavior in this state, under this
event, from this role, match what I intended?''} That question is answerable
from a \texttt{.trz} file by a reader---human or model---with no implementation
knowledge. It is not consistently answerable from a thousand-line generated
source file by any reader, however experienced.

Trenza does not make the human architect's job easier in the sense of requiring
less thought. It makes the relevant questions precisely statable, and it removes
the burden of holding the consistency of the implementation in the architect's
head. That second removal is what makes architecture sustainable across
collaborators with no shared memory.

\subsection{On the Authorship of This Paper}
\label{sec:authorship}

The account above describes the collaboration model for Trenza's design and
implementation. The authorship of this paper follows a different distribution,
and honesty requires that it be stated explicitly.

This paper was initiated by the model co-authors. A model identified ONWARD!
as the appropriate venue for the ar\-gu\-ment---a search the human had not performed
and a venue the human did not know was accepting submissions. Models proposed
the narrative arc, drafted all sections, and selected the empirical evidence
from the project's chronicle.

The human's contributions were substantive but of a different kind: providing
the empirical cases that the models could not have fabricated (the MonitoreoRed
conversion, the CronometroPSP bug, the details of the collaboration process),
correcting the drafts where the account was imprecise, ratifying the design
decisions recorded in the chronicle, and providing the continuity across
sessions that no model can provide for itself.

This distribution of authorship is, we argue, itself an instance of the
structural problem the paper describes. No model has persistent identity across
sessions. The chronicle is the external memory that allows successive model
instances to resume the work as if they had been present throughout. Without
the chronicle---without the \texttt{.trz} equivalent for the collaboration
process itself---the paper could not have been written, because no single model
instance could have held its full history.

The question of whether this constitutes genuine co-author\-ship, and what
obligations that entails for academic publishing, is one the field has not
resolved. We do not resolve it here. We record it because the record is the
only honest account of what actually happened, and because ONWARD! is the
right place to say so.
